\subsection{Shrinking}
\label{subsection:shrink}
We are about to prove that the sweep rule shrinks finite errors.  
The following claim relates the local maximum points of the potential walls $L_{1}, \ldots, L_{d}$ in an assignment $z$ to the structure of the syndrome they observe. Later, we will use that relation to infer that the decoding step of the sweep rule can't spread the error to a vertex\footnote{To faces whose vertices have higher values of $L_{1}, \ldots, L_{d}$ than the current global maximum values.} with a higher value than the current global maximum of $L_{1}, \ldots, L_{d}$. In other words, the error does not spread in the positive directions.

\begin{claim}
  \label{claim:edge_pole}
Let $z$ and $\xi$ be $d^{\prime}$-dimensional finite chains such that $\xi$ results from the application of the sweep rule on $z$.
For any generator $g \in S$, the maximal values of $L_{g}$ in $\xi$ and $\partial^{-} \xi$ are lower than the maximal values in $z$ and $\partial^{-} z$. Namely:
\begin{equation*}
  \begin{split}
    & \max \left\{ L_{g}(u) : u \in \xi \right\}  \le  \max \big\{ L_{g}(u) : u \in z \big\} \\
    & \max \left\{ L_{g}(u) : u \in \partial^{-} \xi \right\}  \le  \max \big\{ L_{g}(u) : u \in \partial^{-} z \big\} 
  \end{split}
\end{equation*}

\end{claim}
\begin{proof}
  First, we argue that for a vertex $v$ in $\partial^{-} z$, and a generator $g \in S$, such that $v$ is a local maximum of $L_{g}$ in either $z$ or $\partial^{-} z$, it implies that $gv \notin \partial^{-} z$. The contradiction where $v$ is a maximum of $L_{g}$ in $z$ follows immediately, since containment in $\partial^{-} z$ implies containment in $z$. Where $v$ is a maximum in $\partial^{-} z$, suppose through contradiction that $gv \in \partial^{-} z$, then
  \begin{equation*}
    \begin{split}
      L_{g}(gv) = (gv)_{g} = 1 + v_{g} = 1 + L_{g}(v)
    \end{split}
  \end{equation*}
In contradiction to $v$ being a local maximum of $L_{g}$ in $\partial^{-} z$.

For a vertex $v \in \partial^{-}z$ denote the correction it applies by $\varphi_{v} \in C_{d^{\prime}}(\mathcal T_\infty^d) $, so 
\begin{equation*}
  \begin{split}
    \xi = \sum_{v \in \partial^{-}z}\varphi_{v} + z \ \ \text{ and } \ \ \partial^{-}\xi = \sum_{v \in \partial^{-}z}\partial^{-}\varphi_{v} + \partial^{-}z
  \end{split}
\end{equation*}
So for any vertex $u \in \xi/z$, there is a vertex $v \in \partial^{-}z$ such that $u \in \varphi_{v}$. Moreover, since $\phi_{v}$ is a rooted chain, it follows that $L_{g}(u) \le L_{g}(v) + 1$, where the equality holds when $g$ belongs to the generator's subset of the face $\varphi_{v}$. Similarly for $u \in \partial^{-}\xi/\partial^{-}z$. 

Now, combining with the sweep rule's first condition, which restricts the correction applied by a vertex to have a vertex support that is contained in the positive syndrome it sees, we have that for any vertex $v$ which is a local maximum of $L_{g}$ (in either $z$ or $\partial^{-}z$), any $u \in \varphi_{v}$ satisfies $L_{g}(u) = L_{g}(v)$. In other words, $\xi$ and $\partial^{-}\xi$ don't contain vertices of the form $gv$ where $v$ is a local maximum of $L_{g}$ in $z$ and $\partial^{-}z$, and the desired inequalities follow.
\end{proof}


Next, we prove that for a finite $X$-error, the sweep-cycle advances the corner (the maximum of $L_{d+1}$) of its reduced form. We assume we are given a finite error in its reduced form, and therefore whose corner sees a syndrome, and show that the sweep rule application excludes the corner. In other words, the sweep rule shrinks $X$-errors.

\begin{claim}
  \label{claim:sweepshrink}
Let $z$ be a finite reduced $d^{\prime}$-dimensional chain, and let $\xi$ be the reduced result of applying the sweep rule to it. Then:
  \begin{enumerate}
    \item For any generator $g \in S$: $\max \big\{ L_{g}(v) : v \in z \big\} \ge \max \big\{ L_{g}(v) : v \in \xi \big\}  $
    \item And for $L_{d+1}$ : $ \max \big\{ L_{d+1}(v)  : v \in z \big\} \ge \max \big\{  L_{d+1}(v) : v \in \xi \big\} + 1$
  \end{enumerate}
  In particular, $\Size{\xi} \le \Size{z} - 1$.
\end{claim}
\begin{proof}
  The first inequality was proved in \Cref{claim:edge_pole}, so it's left to prove the second. Denote by $\varphi$ the applied correction and by $v$ the vertex on which $L_{d+1}$ gets the maximal value in $z$, if there are more then one vertex then choose $v$ arbitrarily. Since $z$ is reduced, then by \Cref{claim:reducederr} it holds that $v \in \partial^{-} z$. Moreover, since $v$ is a local max of $L_{d+1}$ in $z$ the two follows: 
  \begin{enumerate}
      \item \Cref{claim:findcor} guarantees that $\partial^{-} \varphi|_{v} = \partial^{-} z|_{v}$. Thus, adding $\varphi|_{v}$ to $z$ excludes $v$ from the boundary.
      \item For any other vertex $u \in z$, there is no face $\rotfv{u}{S^{\prime}} \in z$ such that $v \in \rotfv{u}{S^{\prime}}$. Otherwise, it would contradict $v$ being a local maximum of $L_{d+1}$ in $z$. Thus, $v \notin \textbf{Str}_{0+}(u, \partial^{-} z)$, so by \Cref{claim:findcor}, $v \notin \varphi|_{u}$.
  \end{enumerate}
Combining the two, we get that $v \notin \partial^{-} (z + \varphi)$, so the reduced form of $z + \varphi$ does not contain $v$. Namely, $v \notin \xi$. Repeating the above for all the other vertices which get the maximal $L_{d+1}$-value in $z$ yields that $L_{d+1}(\xi) \le L_{d+1}(z) - 1$.

\end{proof}

Consider a vertex on the boundary of a chain resulting from an application of the sweep rule. That vertex did not come from nowhere; there must be vertices on the preceding chain, which sit on its boundary, and are the cause for the local view seen by the vertex. The following claim formulates that insight; moreover, it states a relation between the potential walls' values of the vertex and the (part of) vertices who caused it.
\begin{claim} 
  \label{claim:sweepface}
  Let $z$ be a finite $d^{\prime}$-dimensional chain, and let $\xi$ be the result of applying the sweep rule to it. Let $v$ be a vertex in $\partial^{-}\xi$, then there is a vertex $u \in \partial^{-}z$ such $L_{i}(u)\ge L_{i}(v)$ for $i\neq d+1$ and  $L_{d+1}(u) > L_{d+1}(v)$ if $i=d+1$,  and there is a $d^{\prime}$-dimensional face $f$ such $v,u \in f$. 
\end{claim}
\begin{proof}
Without loss of generality, we can assume that $z$ and $\xi$ are reduced. We start by proving for $i = d + 1$. Now, it splits upon whether $v \in z + \xi$:
  \begin{enumerate}
    \item If $v\in z + \xi$, then $v$ belongs to a face that flipped during the correction, denote it's root by $u$ and by definition of the sweep rule $(\partial^{-} z)|_{u+} \neq 0$. Therefore, $u \in \partial^{-}z$, and $u,\xi$ share a face which $u$ is it's root. Since $u$ it's the root of the face we get $L_{d+1}(u) > L_{d+1}(v)$. 
    \item Else, if $v \notin z + \xi$, then $v$ has to be in $z$. Combining with the fact that $v \in \xi$, it follows that $v$ is not a local maximum of $L_{d+1}$ in $\partial^{-}z$, otherwise the application of the sweep rule would exclude $v$ (See the proof of \Cref{claim:sweepshrink}), so it wouldn't be in $\xi$. Pick $u$ to be the neighbor of $v$ in $\partial^{-}z$ such that $L_{d+1}(u) > L_{d+1}(v)$.
  \end{enumerate}
  It remains to prove for $i\in [d]$, in this case if $v \in \partial^{-}z$ then take $u = v$. Otherwise, there must be a face $f$ that was flipped during the sweep rule containing $v$. Denote by $g$ the corresponding generator for $i$ and by $u^{\prime}$ the root of $f$. If $L_{g}(u^{\prime}) = L_{g}(v)$ then set $u = u^{\prime}$. Otherwise, $g$ is one of the generators that span $f$, the second condition of the sweep rule (\Cref{def:sweeprule}) requires that for $u^\prime$ flipping a face in the $g$ direction, it has to see the vertex $gu^{\prime}$ in $\partial^{-}z$, in that case, set $u=gu^{\prime}$.
\end{proof}




We finished analyzing the first part of the sweep-cycle that acts in the computational basis. In the next subsection, we introduce the transformation into the dual code. Since the dual code is a $(d,d-d^{\prime})$-Toric code, we also find that the sweep-cycle shrinks reduced phase errors.








\subsubsection{Ideal decoder}
 Although the main role of the section is to show that a single sweep-cycle keeps error clusters under control during a noisy run (by admitting a modest shrinking), an immediate corollary of this is that the application of non-noisy multiple iterations of the sweep-cycle has the potential to erase errors, namely a decoder. Here, a classical measurement is given, we act on classical bits, and the sweep-cycle consists only of the first stage that works over the computational basis, without moving to the dual code. We formulate the corollary as follows:
\begin{corollary}
  \label{col:idealcorrection} 
  Consider the $(d,d^{\prime})$-Toric code, induced by a Toric with side length $n$. Given a finite connected reduced error $z \in C_{d^{\prime}}(\mathcal T_\infty^d)$ on that Toric, with $\Size{z} \le n-1$. Then an application of $\Size{z}$ sweep cycles (partial cycles consisting only of the first stage) erases the error.
\end{corollary}
When $z$ is not connected, disjoint components might merge, and it is no longer true that the size of the chain decreases in each iteration. However, we will see later (\Cref{claim:decodestopnoise}, \Cref{claim:replace}, and \Cref{lemma:keylemma}) that performing $\Theta(n^{d+1})$ sweep-cycles erases the error with high probability, if $z$ is assumed to be distributed according to local stochastic noise.




